\section{Conclusão}
\label{sec:conclusao}

Este artigo descreveu o projeto, a implementação e a validação do Laboratório
404, um jogo 3D educacional/ficcional no qual uma IA local participa do mundo
sem controlá-lo. O objetivo foi atingido no nível de protótipo: as seis provas
de conceito foram concluídas e demonstráveis, a suíte automatizada cobre 210 a
211 verificações do motor e 16 testes do adapter, todas aprovadas, e o jogo
roda a 60~FPS (limitados por \emph{vsync}) em GPU integrada, mantendo conversa
local com latência entre 12,9~s e 22,0~s.

As contribuições práticas são o padrão de \emph{IA diegética com capacidade
limitada} --- processo separado, contrato estruturado e sanitização em duas
camadas ---, o processo incremental orientado a especificação com evidências
automatizadas e visuais, e os artefatos públicos que permitem reproduzir e
auditar cada afirmação quantitativa deste texto \citep{lab404repo,lab404video}.

As limitações são conhecidas e declaradas: ausência de estudo com usuários,
medições de latência em pequena amostra, validação manual de hardware
pendente, empacotamento independente ainda não gerado e modelo-alvo menor não
executado. Como trabalhos futuros, pretende-se: (i) executar e comparar o
\texttt{llama3.2:3b} com o modelo usado aqui; (ii) concluir o empacotamento e a
demonstração sem intervenção do desenvolvedor (FR-033); (iii) realizar o
playthrough humano de ponta a ponta e, depois, um estudo controlado com
usuários em contexto educacional; e (iv) ampliar o conteúdo do vertical slice
mantendo a mesma disciplina de evidências.
